% Teil II: durchgerechneter Weg, Teile a--d.

\subsection*{Gegeben und gesucht}
Zylinderfoermige Walze (Masse $m$) ruht in einer Rinne; linkes Brett
unter $\alpha_1$, rechtes unter $\alpha_2$ zur Horizontalen.
Gesucht: (a) Freikörperbild mit allen Kräften, (b) Betrag der
Gesamtkraft, (c) Kräfte des linken und rechten Brettes als Funktion von
$\alpha_1,\alpha_2$ und die Gegenkräfte, (d) Winkel für
``vierfache Gewichtskraft'' pro Brett.

\subsection*{Ansatz}
Statik: $\sum\vec F=0$; Krafte in $x$- und $y$-Komponenten zerlegen,
Normalkräfte senkrecht zu den Brettern.

\subsection*{Rechnung}
\textbf{a) Kräfte.} Auf die Walze wirken: die Gewichtskraft
$\vec F_G=m\vec g$ senkrecht nach unten, die Normalkraft $\vec N_1$ des
linken sowie $\vec N_2$ des rechten Brettes (jeweils senkrecht zur
Brettfläche); zusammen ergeben sie null.

\textbf{b) Gesamtkraft.} Wegen Ruhe ist
\begin{equation}
  \sum\vec F=\vec F_G+\vec N_1+\vec N_2=\vec 0
  \;\Rightarrow\;|\vec F_{\mathrm{ges}}|=0.
\end{equation}

\textbf{c) Normalkräfte.} In horizontale und vertikale Richtung zerlegen
(mit $\alpha_1$ (linkes Brett nach links oben) und $\alpha_2$
(rechtes Brett nach rechts oben)):
\begin{align}
  x:\quad -N_1\sin\alpha_1+N_2\sin\alpha_2&=0,\\
  y:\quad N_1\cos\alpha_1+N_2\cos\alpha_2&=mg.
\end{align}
Aus der ersten Gleichung $N_2\sin\alpha_2=N_1\sin\alpha_1$;
einsetzen in die zweite (bzw. symmetrisch) liefert:
\begin{equation}
  \boxed{\;N_1=\frac{mg\,\sin\alpha_2}{\sin(\alpha_1+\alpha_2)},\qquad
  N_2=\frac{mg\,\sin\alpha_1}{\sin(\alpha_1+\alpha_2)}\;}.
\end{equation}
Die Gegenkräfte der Walze auf die Bretter sind
$\vec F'_1=-\vec N_1$ und $\vec F'_2=-\vec N_2$.

\textbf{d) Vierfache Gewichtskraft.} Die Bedingung $N_1=N_2=4mg$
kombiniert mit der $x$-Gleichung $N_1\sin\alpha_1=N_2\sin\alpha_2$
ergibt $\sin\alpha_1=\sin\alpha_2$, also (gleiche, spitze Winkel)
$\alpha_1=\alpha_2=\alpha$. Aus der $y$-Gleichung mit $N=4mg$:
\begin{equation}
  4mg\cos\alpha+4mg\cos\alpha=mg
  \;\Rightarrow\;\cos\alpha=\frac18,
  \;\Rightarrow\;\alpha=\arccos\tfrac18\approx\qty{82.8}{\degree}.
\end{equation}

\begin{ergebnis}
(a) Gewicht plus zwei Normalkräfte; (b) Gesamtkraft null; (c)
$N_1=\dfrac{mg\sin\alpha_2}{\sin(\alpha_1+\alpha_2)}$,
$N_2=\dfrac{mg\sin\alpha_1}{\sin(\alpha_1+\alpha_2)}$, Gegenkräfte $-N_1,-N_2$;
(d) für $N_1=N_2=4mg$: $\alpha_1=\alpha_2=\arccos(1/8)\approx\qty{82.8}{\degree}$.
\end{ergebnis}

\begin{kontrolle}
\textbf{Sonderfall:} Für gleiche Winkel $\alpha_1=\alpha_2=45^\circ$ ist
$\sin(\alpha_1+\alpha_2)=1$ und
$N_1=N_2=mg\,\sin45^\circ=mg/\sqrt2\approx0{,}71\,mg$. Ein symmetrischer
45$^\circ$-Keil ruft an jedem Brett eine Normalkraft von $0{,}71\,mg$
hervor; die Summe der Vertikalkomponenten ($2\cdot0{,}71\,mg\cos45^\circ$)
gleicht $mg$ -- plausibel. \checkmark

\textbf{Einheiten:} $N_1,N_2$ in \unit{\newton}; Argumente $mg\sin\alpha_2$
[\unit{\newton}], $\sin(\alpha_1+\alpha_2)$ dimensionslos. \checkmark
\end{kontrolle}

\begin{merksatz}
Im statischen Gleichgewicht ist die Vektorsumme aller Kräfte null. Zwei
schraege Auflager zerlegt man in Normalkräfte mit einer
Komponentenzerlegung; die Kraft pro Brett hängt von beiden Winkeln ab und
kann für steile Rinnen die Gewichtskraft um ein Vielfaches übersteigen.
\end{merksatz}
